\section{Kết luận}
\label{sec:conclusion}

Đồ án đã hoàn thành việc nghiên cứu, cài đặt và đánh giá thuật toán FP-Max cho bài toán khai thác tập phổ biến tối đại. Về lý thuyết, báo cáo trình bày các định nghĩa cốt lõi của FIM, tính chất Apriori, cấu trúc FP-tree làm nền, và phần trọng tâm về FP-Max: danh sách MFI, phép tỉa theo siêu tập, shortcut single-path và phân tích độ phức tạp. Hai ví dụ tay minh hoạ FP-Max trên cả trường hợp cơ sở (cập nhật danh sách MFI) lẫn tình huống đặc biệt (phép tỉa bỏ qua nguyên một nhánh).

Về cài đặt, nhóm xây dựng package Julia \texttt{FrequentItemsetMining} gồm FP-Growth (bản cơ bản và bản tối ưu mã hoá số nguyên), một baseline naive-maximal khai thác toàn bộ tập phổ biến rồi lọc tối đại, và FP-Max khai thác trực tiếp có danh sách MFI và tỉa theo siêu tập. Output được giữ tương thích với định dạng SPMF để thuận tiện kiểm thử và đối chiếu.

Về thực nghiệm, kết quả correctness đạt \texttt{match\_ratio = 1.0} và \texttt{support\_mismatch = 0} so với SPMF FPMax trên mọi cấu hình benchmark được kiểm tra. Điểm mạnh quyết định của FP-Max trực tiếp là chạy được trong những vùng mà baseline naive-maximal không khả thi --- toàn bộ Accidents và các ngưỡng thấp trên dữ liệu dày --- nhờ chỉ giữ biểu diễn tối đại và tỉa không gian tìm kiếm. Trên dữ liệu dày, FP-Max còn nhanh hơn baseline nhiều lần, ví dụ khoảng 8.68 lần trên Chess tại minsup $0.8$ và 4.84 lần trên Mushrooms tại minsup $0.25$. Tuy vậy, trên dữ liệu thưa (Retail, T10I4D100K) FP-Max chậm hơn baseline vì phép tỉa ít tác dụng còn chi phí cài đặt cao, và nhìn chung vẫn chậm hơn SPMF.

Phần ứng dụng trên Groceries dùng bộ máy khai thác đầy đủ tập phổ biến (vì sinh luật cần support của mọi itemset) để tạo các luật kết hợp có ý nghĩa thực tế, chẳng hạn các luật liên quan tới \textit{root vegetables}, \textit{other vegetables}, \textit{whole milk} và \textit{yogurt}. Các luật này có thể hỗ trợ gợi ý sản phẩm, bố trí kệ hàng và thiết kế combo.

Các hạn chế còn lại của đồ án gồm: cài đặt FP-Max hiện dùng FP-tree với node \texttt{String} và tính lại support cho từng MFI bằng một lần quét toàn cơ sở dữ liệu, khiến nó chậm trên dữ liệu thưa; baseline naive-maximal phải skip nhiều cấu hình do nguy cơ bùng nổ. Trong tương lai, nhóm có thể chuyển FP-Max sang cấu trúc cây mã hoá số nguyên như bộ máy tối ưu, mang support theo trong quá trình khai thác để bỏ bước tính lại, thay danh sách MFI tuyến tính bằng một MFI-tree cho phép kiểm tra siêu tập dưới tuyến tính, và song song hoá các nhánh conditional tree độc lập.
